\documentclass[a4paper,11pt]{article}

%%% Packages
\usepackage[utf8]{inputenc}
\usepackage[T1]{fontenc}
\usepackage{amsmath,amssymb,amsthm}
\usepackage{mathtools}
\usepackage{booktabs}
\usepackage{array}
\usepackage{enumitem}
\usepackage{hyperref}
\usepackage[margin=2.5cm]{geometry}
\usepackage{xcolor}

\hypersetup{colorlinks=true,linkcolor=blue!45!black,citecolor=blue!45!black,urlcolor=blue!45!black}

%%% Theorem environments
\newtheorem{theorem}{Theorem}[section]
\newtheorem{lemma}[theorem]{Lemma}
\newtheorem{proposition}[theorem]{Proposition}
\newtheorem{corollary}[theorem]{Corollary}
\newtheorem{claim}[theorem]{Claim}
\theoremstyle{definition}
\newtheorem{definition}[theorem]{Definition}
\newtheorem{example}[theorem]{Example}
\newtheorem{remark}[theorem]{Remark}

%%% Notation macros
\newcommand{\ALCI}{\ensuremath{\mathcal{ALCI}}}
\newcommand{\ALCIRCC}[1]{\ensuremath{\mathcal{ALCI}_{\mathit{RCC#1}}}}
\newcommand{\DR}{\ensuremath{\mathsf{DR}}}
\newcommand{\PO}{\ensuremath{\mathsf{PO}}}
\newcommand{\EQ}{\ensuremath{\mathsf{EQ}}}
\newcommand{\PP}{\ensuremath{\mathsf{PP}}}
\newcommand{\PPI}{\ensuremath{\mathsf{PPI}}}
\newcommand{\inv}{\ensuremath{\mathsf{inv}}}
\newcommand{\comp}{\ensuremath{\mathsf{comp}}}
\newcommand{\cl}{\ensuremath{\mathsf{cl}}}
\newcommand{\tp}{\ensuremath{\mathsf{tp}}}
\newcommand{\Tp}{\ensuremath{\mathsf{Tp}}}
\newcommand{\NR}{\ensuremath{\mathsf{NR}}}
\newcommand{\Tri}{\ensuremath{\mathsf{Tri}}}
\newcommand{\TNbr}{\ensuremath{\mathsf{TNbr}}}
\newcommand{\EXPTIME}{\textsc{ExpTime}}
\newcommand{\Safe}{\ensuremath{\mathsf{Safe}}}
\newcommand{\interp}{\ensuremath{\mathcal{I}}}
\newcommand{\DeltaI}{\ensuremath{\Delta^{\mathcal{I}}}}
\newcommand{\GG}{\ensuremath{\mathcal{G}}}
\newcommand{\LL}{\ensuremath{\mathcal{L}}}
\newcommand{\EE}{\ensuremath{\mathcal{E}}}
\newcommand{\NN}{\ensuremath{\mathcal{N}}}

\setlength{\parindent}{0pt}
\setlength{\parskip}{0.55em}

\title{A Tableau Calculus for $\ALCIRCC{5}$\\
  with Tri-Neighborhood Blocking\\[4pt]
  \large Second revision, April 2026}

\author{Claude (Anthropic) \and Michael Wessel}

\date{April 2026}

\begin{document}

\maketitle

\begin{abstract}
We present a tableau-based decision procedure for concept
satisfiability in $\ALCIRCC{5}$, the description logic $\ALCI$
extended with RCC5 composition-based role interactions over a
complete relational structure. The key innovation is
\emph{Tri-neighborhood blocking}: a node $x$ is blocked by an
earlier node $y$ only when they share the same concept label,
the same set of abstract triangle types---tuples
$(\tau_1, R_{12}, \tau_2, R_{23}, \tau_3, R_{13})$ recording
Hintikka types and RCC5 relations but no node identities---and,
for each (relation, type) pair, the same distribution of
triangle-type sets among their respective neighbors.
We prove termination (the neighborhood signature is drawn
from a finite universe), soundness (the model is constructed
by triangle-type-filtered unraveling, with global consistency
guaranteed by the full tractability of RCC5), and completeness
(a model guides all nondeterministic choices).
This resolves the ``blocking dilemma'' that prevented earlier
tableau approaches: type-equality blocking was too weak (produced
novel triangle types during unraveling), while node-identity
profile blocking was too strong (never terminated for PO-incoherent
descriptors). Tri-neighborhood blocking achieves both properties
simultaneously, and the strengthened neighborhood condition
ensures that the copy is faithful from every participating
node's perspective.
\end{abstract}


%% ============================================================
\section{Introduction}\label{sec:intro}

The description logic $\ALCIRCC{5}$ extends $\ALCI$ with a
role box defined by the RCC5 composition table: the five base
relations $\{\DR, \PO, \EQ, \PP, \PPI\}$ act as role names,
and every pair of domain elements is related by exactly one
base relation, subject to global composition
consistency~\cite{Wessel2002}.  Models are thus complete
graphs with composition-consistent edge labels---a fundamental
departure from tree-shaped $\ALCI$ models.

The decidability of $\ALCIRCC{5}$ has remained open since
Wessel~\cite{Wessel2002} established decidability of
$\ALCI_{\mathit{RCC3}}$ and proved $\ALCIRCC{5}$ PSPACE-hard.
Several approaches have been attempted: a quasimodel method
with type elimination~\cite{CompanionPaper}, a contextual
tableau with local states~\cite{GPTTableau}, direct model
construction~\cite{ClosingGap}, and a two-tier
quotient~\cite{TwoTier} that proves decidability for the
\emph{PO-coherent} fragment.

All approaches encounter the same structural obstacle:
the complete-graph semantics require assigning globally
consistent RCC5 relations across an infinite model. The
\emph{patchwork property} of RCC5~\cite{RenzNebel1999}---that
path-consistent atomic networks are globally consistent---bridges
local to global consistency, but the question is whether the
finite information in a tableau completion graph suffices to
guarantee path-consistency of the infinite unraveled structure.

\subsection*{The blocking dilemma}

In a standard DL tableau, blocking ensures termination:
a blocked node $x$ copies the witness structure of its blocker
$y$ during model extraction.  In complete-graph models, this
copying (``unraveling'') must preserve RCC5 composition
consistency across \emph{all} triples, not just the
parent-child edge.

This creates a dilemma, identified in~\cite{TwoTier}:
\begin{itemize}[nosep]
  \item \textbf{Type-equality blocking} (block when
    $\LL(x) = \LL(y)$) always terminates but may produce
    \emph{novel triangle types} during unraveling---triples
    $(\tau_1, R, \tau_2, S, \tau_3, R')$ not witnessed in
    the completion graph---potentially violating
    path-consistency.
  \item \textbf{Node-identity profile blocking} (block when
    $x$ and $y$ have identical relational profiles to all
    other nodes) always supports correct unraveling but may
    \emph{never terminate}: for PO-incoherent descriptors,
    the ``sliding PO diagonal'' gives each node a unique
    profile.
\end{itemize}

\subsection*{Tri-neighborhood blocking}

We resolve this dilemma by observing that the non-termination
argument uses \emph{concrete} node identities, while the
correctness argument only needs \emph{abstract} relational
patterns.  An \emph{abstract triangle type} is a tuple
$(\tau_1, R_{12}, \tau_2, R_{23}, \tau_3, R_{13})$ recording
three Hintikka types and three pairwise RCC5 relations---no
node identities.  The \emph{triangle-type set} of a node $x$
is the set of all abstract triangle types that $x$ participates
in.

Computational verification on the PO-incoherent counterexample
(a 24-element PP-chain with 12 PO-witnesses) shows that abstract
triangle-type sets \emph{stabilize}: all interior $\tau_A$ nodes
share exactly the same 68 abstract triangle types, despite having
different concrete neighbors.  The stabilization occurs because
the abstract triangle-type set depends on the \emph{menu} of
available (type, relation) pairs among neighbors, not on the
\emph{count} of neighbors of each kind.

We strengthen this to \emph{Tri-neighborhood blocking}:
in addition to requiring $\LL(x) = \LL(y)$ and
$\Tri(x) = \Tri(y)$, we require that for each
(relation, type) pair, the set of $\Tri$-values among
the respective neighbors of $x$ and $y$ also matches.
This ensures the copy is faithful not just from $x/y$'s
perspective, but from \emph{every neighboring node's
perspective}.  Computational verification confirms that this
stronger condition also stabilizes (at depth $k = 3$, one
step later than basic $\Tri$-matching at $k = 2$).

Tri-neighborhood blocking achieves:
\begin{enumerate}[nosep]
  \item \textbf{Termination}: the neighborhood signature is
    drawn from a finite universe, and stabilization ensures
    matching.
  \item \textbf{Correct unraveling}: copying $y$'s witness
    structure for $x$ produces only triangle types already in
    the completion graph's set $T$, verified from every
    participating node's perspective.
  \item \textbf{Global consistency}: $T$-closure plus full
    RCC5 tractability~\cite{RenzNebel1999,Renz1999} yields a
    globally consistent atomic network.
\end{enumerate}


%% ============================================================
\section{Preliminaries}\label{sec:prelim}

\subsection{The logic $\ALCIRCC{5}$}

Concepts are built from concept names, the RCC5 base relations
$R \in \{\DR, \PO, \EQ, \PP, \PPI\}$ as role names, and the
constructors $\sqcap$, $\sqcup$, $\neg$, $\exists R.C$,
$\forall R.C$ in negation normal form (NNF).  We adopt
\textbf{strong $\EQ$ semantics}: $\rho(d, e) = \EQ$ iff
$d = e$.

An interpretation $\interp = (\DeltaI, \cdot^{\interp})$
assigns a total function
$\rho^{\interp}: \DeltaI \times \DeltaI \to
\{\DR, \PO, \EQ, \PP, \PPI\}$ subject to: (i)~$\rho(d,d) = \EQ$;
(ii)~$\rho(d,e) = \inv(\rho(e,d))$;
(iii)~composition consistency:
$\rho(d_1, d_3) \in \comp(\rho(d_1, d_2), \rho(d_2, d_3))$
for all $d_1, d_2, d_3$; and the usual concept-interpretation
clauses for $\ALCI$.

\subsection{Types, closure, composition table}

Given $C_0$ in NNF, $\cl(C_0)$ is the Fischer--Ladner closure.
A \emph{Hintikka type} $\tau \subseteq \cl(C_0)$ is a maximal
propositionally consistent subset.  Write $\Tp(C_0)$ for the
set of types; $|\Tp(C_0)| \le 2^{|\cl(C_0)|}$.

Write $\NR^{-} = \{\DR, \PO, \PP, \PPI\}$ for the base
relations between distinct elements ($\EQ$ excluded).

The composition table
$\comp: (\NR^{-})^2 \to 2^{\NR^{-}}$ is:

\begin{center}
\renewcommand{\arraystretch}{1.15}
\begin{tabular}{c|cccc}
$\comp(R, S)$ & $\DR$ & $\PO$ & $\PP$ & $\PPI$ \\
\hline
$\DR$ & \{all\} & \{$\DR$,$\PO$,$\PP$\} &
  \{$\DR$,$\PO$,$\PP$\} & \{$\DR$\} \\
$\PO$ & \{$\DR$,$\PO$,$\PPI$\} & \{all\} &
  \{$\PO$,$\PP$\} & \{$\DR$,$\PO$,$\PPI$\} \\
$\PP$ & \{$\DR$\} & \{$\DR$,$\PO$,$\PP$\} &
  \{$\PP$\} & \{all\} \\
$\PPI$ & \{$\DR$,$\PO$,$\PPI$\} & \{$\PO$,$\PPI$\} &
  \{$\PO$,$\PP$,$\PPI$\} & \{$\PPI$\} \\
\end{tabular}
\end{center}

\subsection{Type-safety}

\begin{definition}[Type-safe relations]\label{def:safe}
For Hintikka types $\tau_1, \tau_2$, define
\[
\Safe(\tau_1, \tau_2) = \bigl\{R \in \NR^{-} :
\forall\, (\forall R.C) \in \tau_1,\ C \in \tau_2;\;
\forall\, (\forall\inv(R).C) \in \tau_2,\ C \in \tau_1
\bigr\}.
\]
In any model, $\rho(d, e) \in \Safe(\tp(d), \tp(e))$ for
distinct $d, e$.
\end{definition}

\subsection{Full tractability of RCC5}

\begin{theorem}[Renz \& Nebel~\protect\cite{RenzNebel1999},
Renz~\protect\cite{Renz1999}]\label{thm:patchwork}
Every path-consistent disjunctive RCC5 constraint network
over $\{\DR, \PO, \EQ, \PP, \PPI\}$ is satisfiable.
Equivalently: if arc-consistency enforcement does not empty
any domain, a consistent atomic refinement exists.\qed
\end{theorem}

This is the \emph{patchwork property}: local consistency
(path-consistency) implies global consistency.


%% ============================================================
\section{The Tableau Calculus}\label{sec:tableau}

\subsection{Completion graphs}

\begin{definition}[Completion graph]\label{def:cg}
A \emph{completion graph} for $C_0$ is a tuple
$\GG = (V, \LL, \EE, \prec)$ where:
\begin{itemize}[nosep]
  \item $V$ is a finite non-empty set of nodes.
  \item $\LL: V \to 2^{\cl(C_0)}$ is a labelling function.
  \item $\EE: \{(x,y) \in V^2 \mid x \neq y\}
    \to \NR^{-}$ assigns a base relation to every ordered pair
    of distinct nodes, with $\EE(y,x) = \inv(\EE(x,y))$.
  \item $\prec$ is a strict total order on $V$ (creation order).
\end{itemize}
The \emph{initial} completion graph has a single node $x_0$
with $\LL(x_0) = \{C_0\}$.
\end{definition}

\subsection{Triangle types}\label{sec:tri-types}

\begin{definition}[Abstract triangle type]\label{def:tri}
An \emph{abstract triangle type} is a tuple
$(\tau_1, R_{12}, \tau_2, R_{23}, \tau_3, R_{13})$ where
$\tau_1, \tau_2, \tau_3 \in \Tp(C_0)$ and
$R_{12}, R_{23}, R_{13} \in \NR^{-}$ with
$R_{13} \in \comp(R_{12}, R_{23})$.
\end{definition}

\begin{definition}[Triangle-type set of a node]\label{def:tri-set}
For a node $x$ in $\GG$, the \emph{triangle-type set} of $x$ is:
\[
\Tri(x) = \bigl\{
  (\LL(x), \EE(x,b), \LL(b), \EE(b,c), \LL(c), \EE(x,c))
  : b, c \in V,\; |\{x,b,c\}| = 3
\bigr\}.
\]
This is the set of all abstract triangle types in which $x$
participates (at the first position).
\end{definition}

\begin{remark}
Since each abstract triangle type uses only Hintikka types
and RCC5 relations---no node identities---the set $\Tri(x)$
is drawn from the finite universe
$\Tp(C_0)^3 \times (\NR^{-})^3$.  Its cardinality is bounded
by $|\Tp|^2 \cdot |\NR^{-}|^2 = 2^{2|\cl(C_0)|} \cdot 16$
(since the first component $\tau_1 = \LL(x)$ is fixed).
\end{remark}

\begin{definition}[Triangle-type set of the graph]\label{def:T}
The \emph{triangle-type set} of $\GG$ is:
\[
T(\GG) = \bigl\{
  (\LL(x), \EE(x,y), \LL(y), \EE(y,z), \LL(z), \EE(x,z))
  : \text{distinct } x, y, z \in V
\bigr\}.
\]
\end{definition}

\subsection{Tri-neighborhood signature}\label{sec:tri-nbr}

\begin{definition}[Tri-neighborhood signature]\label{def:tri-nbr}
For a node $x$ in $\GG$, the \emph{Tri-neighborhood signature}
of $x$ is:
\[
\TNbr(x) = \bigl\{
  \bigl((R, \tau),\; \{\Tri(b) : b \in V,\; b \neq x,\;
    \EE(x,b) = R,\; \LL(b) = \tau\}\bigr)
  : R \in \NR^{-},\; \tau \in \Tp
\bigr\}.
\]
That is, for each (relation, type) pair $(R, \tau)$,
$\TNbr(x)$ records the \emph{set} of distinct $\Tri$-values
occurring among $x$'s $R$-neighbors of type $\tau$.
\end{definition}

\begin{remark}
The Tri-neighborhood signature captures not only which
abstract triangle types $x$ participates in, but also which
abstract triangle types each of $x$'s neighbors participates
in, grouped by how the neighbor relates to $x$.  This is
analogous to \emph{pairwise blocking} in $\ALCI$ tree tableaux
(which matches a node \emph{and} its parent with the blocker
and its parent), generalized to the complete-graph setting:
instead of matching one parent edge, we match \emph{all}
neighbor edges grouped by (relation, type).
\end{remark}

\subsection{Blocking}\label{sec:blocking}

\begin{definition}[Tri-neighborhood blocking]\label{def:blocking}
A node $x \in V$ is \emph{blocked} if there exists $y \prec x$
such that $y$ is not blocked and:
\begin{enumerate}[nosep,label=(\roman*)]
  \item $\LL(x) = \LL(y)$ (same concept label),
  \item $\Tri(x) = \Tri(y)$ (same triangle-type set), and
  \item $\TNbr(x) = \TNbr(y)$ (same Tri-neighborhood signature).
\end{enumerate}
We call $y$ a \emph{blocker} for $x$.  A node is \emph{active}
if it is not blocked.
\end{definition}

\begin{remark}
Conditions~(ii) and~(iii) are the key innovations.  Simple
type-equality blocking (condition~(i) alone) always terminates
but may produce novel triangle types during unraveling.
Condition~(ii) ensures that $x$ and $y$ participate in the
same abstract relational patterns.  Condition~(iii) additionally
ensures that the \emph{neighbors} of $x$ and $y$ also have
matching triangle-type sets: for each (relation, type) pair,
the menu of $\Tri$-values among neighbors is identical.

This strengthening is crucial for the soundness proof: when
we copy $y$'s witness structure for blocked node $x$, we need
the triangle types to be in $T$ not just from $x/y$'s
perspective, but from \emph{every neighbor's perspective}.
Condition~(iii) guarantees this.

Note that conditions~(ii)--(iii) do \emph{not} require $x$ and
$y$ to be related to the same concrete nodes.  They require only
that the \emph{abstract neighborhood patterns} match.  This is
strictly weaker than node-identity profile matching, which is why
the condition still terminates (Section~\ref{sec:termination}).
\end{remark}

\subsection{Expansion rules}\label{sec:rules}

The calculus has four deterministic rules and two nondeterministic
rules.

\medskip
\noindent\textbf{Propositional rules} (deterministic):

\smallskip
\noindent$(\sqcap)$\quad If $D_1 \sqcap D_2 \in \LL(x)$ and
$\{D_1, D_2\} \not\subseteq \LL(x)$: add $D_1$ and $D_2$ to
$\LL(x)$.

\smallskip
\noindent$(\sqcup)$\quad If $D_1 \sqcup D_2 \in \LL(x)$ and
$D_1 \notin \LL(x)$ and $D_2 \notin \LL(x)$:
\emph{nondeterministically} add $D_1$ or $D_2$.

\medskip
\noindent\textbf{Propagation rule} (deterministic):

\smallskip
\noindent$(\forall)$\quad If $\forall R.D \in \LL(x)$ and
$\EE(x,y) = R$ and $D \notin \LL(y)$: add $D$ to $\LL(y)$.

\smallskip
\noindent This handles inverses: $\forall R^-.D$ becomes
$\forall \inv(R).D$ in NNF.

\medskip
\noindent\textbf{Generating rule} (nondeterministic):

\smallskip
\noindent$(\exists)$\quad If $\exists R.D \in \LL(x)$, $x$ is
active, and no $y \in V$ satisfies $\EE(x,y) = R$ and
$D \in \LL(y)$:
\begin{enumerate}[nosep,label=\arabic*.]
  \item Create a fresh node $y$, $\prec$-greatest.
  \item Set $\LL(y) := \{D\}$.
  \item Set $\EE(x,y) := R$, $\EE(y,x) := \inv(R)$.
  \item For each $z \in V \setminus \{x,y\}$,
    \emph{nondeterministically} choose
    $S_z \in \NR^{-}$ and set $\EE(z,y) := S_z$,
    $\EE(y,z) := \inv(S_z)$.
\end{enumerate}

\subsection{Constraint filtering}\label{sec:filtering}

After step~4 of $(\exists)$, we require:

\medskip\noindent\textbf{(CF) Composition consistency.}\;
For every triple of distinct nodes $(u,v,w)$ involving $y$:
\[
\EE(u,w) \in \comp(\EE(u,v), \EE(v,w)).
\]
If no assignment of $\{S_z\}$ satisfies (CF) for all triples
simultaneously, the branch \emph{closes} (clash by composition
failure).

\begin{remark}
Checking (CF) is an RCC5 constraint satisfaction problem.  By
Theorem~\ref{thm:patchwork}, a simultaneous satisfying
assignment exists iff path-consistency enforcement does not
empty any domain.  This can be checked in $O(|V|^3)$.
\end{remark}

\subsection{Rule application strategy}\label{sec:strategy}

\begin{enumerate}[nosep]
  \item Apply $(\sqcap)$, $(\sqcup)$, $(\forall)$ exhaustively
    to all nodes.
  \item Evaluate blocking: $x$ is blocked if $\LL(x) = \LL(y)$,
    $\Tri(x) = \Tri(y)$, and $\TNbr(x) = \TNbr(y)$ for some
    non-blocked $y \prec x$.
  \item Apply one instance of $(\exists)$ for an active node
    with an unsatisfied demand.
  \item Apply $(\forall)$ exhaustively (new edges may trigger
    propagation).
  \item Re-evaluate blocking.
  \item Repeat from step~1 until no rule applies or a clash
    is detected.
\end{enumerate}

\subsection{Clash conditions}

\begin{definition}[Saturated and clash-free]\label{def:open}
$\GG$ is \emph{saturated} if no rule applies. $\GG$ is
\emph{clash-free} if:
\begin{enumerate}[nosep,label=(C\arabic*)]
  \item No $x$ has $\{A, \neg A\} \subseteq \LL(x)$.
  \item Composition consistency (CF) holds for all triples.
\end{enumerate}
A saturated, clash-free graph is \emph{open}.
\end{definition}


%% ============================================================
\section{Termination}\label{sec:termination}

\begin{theorem}[Termination]\label{thm:termination}
Every sequence of rule applications terminates.
\end{theorem}

\begin{proof}
Let $n = |\cl(C_0)|$ and $N_\tau = |\Tp(C_0)| \le 2^n$.

\medskip\noindent\textbf{Finiteness of blocking signatures.}\;
An abstract triangle type is an element of
$\Tp^3 \times (\NR^{-})^3$.  Let $S$ denote the number of
possible $\Tri$-values; $S \le 2^{N_\tau^3 \cdot 64}$.
A Tri-neighborhood signature assigns, to each of the at most
$4 \cdot N_\tau$ pairs $(R, \tau)$, a subset of $\{1, \ldots, S\}$
(a set of $\Tri$-values).  The number of distinct
$\TNbr$-values is at most $(2^S)^{4 N_\tau}$.
The number of distinct $(\LL, \Tri, \TNbr)$ triples is
therefore bounded by
$A(C_0) := N_\tau \cdot S \cdot (2^S)^{4 N_\tau}$.

\medskip\noindent\textbf{Active node bound.}\;
Active nodes have pairwise distinct
$(\LL, \Tri, \TNbr)$ triples (by the blocking
definition).  Hence the number of active nodes is at most
$A(C_0)$.

\medskip\noindent\textbf{Labels grow monotonically.}\;
Rules only add concepts to labels, never remove them.  Each
label changes at most $n$ times.

\medskip\noindent\textbf{Bounded configuration changes.}\;
Adding a new node can add triangle types for existing nodes.
Changing a label (via $\forall$-propagation) can change the
Hintikka-type component of existing triangle types and thus
alter $\Tri$ and $\TNbr$ values.  Since labels only grow,
at each node the number of distinct
$(\LL, \Tri, \TNbr)$ configurations encountered is
bounded by $n$ times the number of possible
$(\Tri, \TNbr)$ configurations---a finite quantity.

\medskip\noindent\textbf{Node creation bound.}\;
Each active node has at most $n$ existential demands. With
at most $A(C_0)$ active nodes, at most $A(C_0) \cdot n$ nodes
are created.  Re-evaluations of blocking after label changes
or node additions may unblock nodes, but each unblocking
corresponds to a configuration change, and the total is bounded.

\medskip\noindent\textbf{Overall.}\;
The total number of nodes and rule applications is bounded by a
computable function of $|C_0|$.  The bound is non-elementary
(due to the tower of exponentials in $A(C_0)$), but finite.
\end{proof}

\begin{remark}[Tighter bounds]\label{rem:bounds}
The bound $A(C_0)$ is astronomically large.  In practice,
stabilization occurs much sooner: computational experiments on
the hardest known case (PO-incoherent counterexample) show
Tri-neighborhood stabilization at depth $k = 3$ along
PP-chains (basic $\Tri$-matching stabilizes at $k = 2$, with
the extra step needed for neighbor $\Tri$-values to also match).
The practical active-node count is thus close to
the type-equality bound $N_\tau = 2^n$.

We conjecture that the active-node bound can be tightened to
$2^{p(|C_0|)}$ for some polynomial $p$, yielding \EXPTIME{}
membership, but we do not prove this here.
\end{remark}


%% ============================================================
\section{Soundness}\label{sec:soundness}

\begin{theorem}[Soundness]\label{thm:soundness}
If the tableau produces an open completion graph $\GG$ for
$C_0$, then $C_0$ is satisfiable in $\ALCIRCC{5}$.
\end{theorem}

The proof constructs a model directly from $\GG$ via tree
unraveling and triangle-type-filtered constraint solving.

\subsection{Tree unraveling}\label{sec:unraveling}

\begin{definition}[Tree unraveling]\label{def:unrav}
Let $V_a \subseteq V$ be the active nodes.  For each blocked
node~$x$, let $\beta(x)$ denote its blocker.  The tree
unraveling $(\Delta, \mathsf{map}, \mathsf{par}, \sigma)$ is
defined inductively:
\begin{enumerate}[nosep]
  \item The root $d_0 \in \Delta$ has $\mathsf{map}(d_0) = x_0$.
  \item For each $d \in \Delta$ with $\mathsf{map}(d) = n$,
    and each existential demand $\exists S.D \in \LL(n)$
    witnessed by node~$w$ (i.e.,
    $\EE(n, w) = S$ and $D \in \LL(w)$): create a child $d'$
    with $\mathsf{par}(d') = d$, $\sigma(d') = S$, and
    \[
    \mathsf{map}(d') = \begin{cases}
      w & \text{if } w \in V_a, \\
      \beta(w) & \text{if } w \text{ is blocked.}
    \end{cases}
    \]
\end{enumerate}
Each element's type is $\tp(d) = \LL(\mathsf{map}(d))$.
\end{definition}

\begin{remark}
When a blocked node $w$ is redirected to its blocker $\beta(w)$,
the child $d'$ maps to $\beta(w)$ and inherits $\beta(w)$'s
label.  Since $\LL(w) = \LL(\beta(w))$ (condition~(i) of
blocking), the type is preserved.  The critical property is
that $\beta(w)$'s existential demands, when recursively
expanded, produce the \emph{same abstract triangle types} as
$w$ would have---guaranteed by conditions~(ii) and~(iii) of
blocking: $\Tri(w) = \Tri(\beta(w))$ and
$\TNbr(w) = \TNbr(\beta(w))$.  The Tri-neighborhood condition
further ensures that each witness $v$ of $\beta(w)$ at relation
$R$ of type $\tau$ has a matching $\Tri$-value to the
corresponding witness of $w$, so the copy is faithful from
every participant's perspective.
\end{remark}

\subsection{The triangle-type-filtered constraint network}
\label{sec:network}

\begin{definition}[Constraint network $\NN_T$]\label{def:network}
Let $T = T(\GG)$ and $P = P(\GG) = \{(\LL(x), \EE(x,y), \LL(y))
: x \neq y \in V\}$.  For each pair of distinct elements
$d_1, d_2 \in \Delta$:

\medskip\noindent\textbf{Parent--child edges:}\;
If $\mathsf{par}(d_2) = d_1$, then $\rho(d_1, d_2) = \sigma(d_2)$
(fixed).

\medskip\noindent\textbf{Non-adjacent pairs:}\;
The initial domain is
$D_0(d_1, d_2) = \{R \in \NR^{-} : (\tp(d_1), R, \tp(d_2)) \in P\}$.

\medskip
The \emph{triangle-filtered constraint network} $\NN_T$ is
obtained by enforcing arc-consistency: iteratively remove any
$R \in D(d_1, d_2)$ for which there exists $d_3$ such that for
some $R_{23} \in D(d_2, d_3)$, no $R_{13} \in D(d_1, d_3)$
satisfies
$(\tp(d_1), R, \tp(d_2), R_{23}, \tp(d_3), R_{13}) \in T$.
\end{definition}

\subsection{The key lemma: $T$-closed solutions exist}
\label{sec:key-lemma}

\begin{definition}[$T$-closed solution]\label{def:tclosed}
A \emph{$T$-closed solution} is an assignment
$\rho: \Delta \times \Delta \to \NR^{-}$ such that:
\begin{enumerate}[nosep,label=(\roman*)]
  \item $\rho(d_1, d_2) \in D_0(d_1, d_2)$ for non-adjacent pairs.
  \item $\rho(d_1, d_2) = \sigma(d_2)$ for parent--child pairs.
  \item For every triple of distinct $d_1, d_2, d_3$:
    $(\tp(d_1), \rho(d_1,d_2), \tp(d_2), \rho(d_2,d_3),
    \tp(d_3), \rho(d_1,d_3)) \in T$.
\end{enumerate}
\end{definition}

\begin{lemma}[Existence of $T$-closed solutions]
\label{lem:tclosed-exist}
Let $\GG$ be an open completion graph with Tri-neighborhood
blocking.  Then the constraint network $\NN_T$ admits a
$T$-closed solution.
\end{lemma}

\begin{proof}
We construct $\rho$ by induction on the tree structure of $\Delta$.
The assignment rule is: for distinct elements $d_1, d_2$ with
$\mathsf{map}(d_1) = n_1$ and $\mathsf{map}(d_2) = n_2$,
set $\rho(d_1, d_2) = \EE(n_1, n_2)$ when $n_1 \neq n_2$
(the case $n_1 = n_2$ is handled separately below).

\medskip\noindent\textbf{Base.}\;
For the root $d_0$ and its children: parent--child edges
are fixed.  For pairs of children $d_i, d_j$ with maps
$n_i = \mathsf{map}(d_i)$, $n_j = \mathsf{map}(d_j)$:
set $\rho(d_i, d_j) = \EE(n_i, n_j)$.  All triples among
$\{d_0, d_i, d_j\}$ map to triples of nodes in $\GG$, so
their triangle types are in $T(\GG)$.

\medskip\noindent\textbf{Cross-subtree edges (distinct maps).}\;
Consider $d_1$ with map $n_1$ and $d_2$ with map $n_2$, where
$n_1 \neq n_2$.  Set $\rho(d_1, d_2) = \EE(n_1, n_2)$.
For any third element $d_3$ with map $n_3$: if
$|\{n_1, n_2, n_3\}| = 3$, then $(n_1, n_2, n_3)$ is a triple
in $\GG$ and the triangle type is in $T(\GG)$.  If $n_3 = n_1$
or $n_3 = n_2$, see the same-map case below.

\medskip\noindent\textbf{Intra-subtree edges via blocked copies.}\;
When $d_2$ is in a subtree rooted at a blocked copy of node $w$
that redirects to $\beta(w)$, the expansion of $\beta(w)$ is
used in place of $w$'s.  Consider a triple $(d_1, d_2, d_3)$
where $d_2$ maps to some $n_2$ in $\beta(w)$'s expansion.

\emph{This is where Tri-neighborhood blocking is essential.}
We need $(\tp(d_1), \rho(d_1, d_2), \tp(d_2), \rho(d_2, d_3),
\tp(d_3), \rho(d_1, d_3)) \in T$.  There are two cases:

\smallskip\noindent\emph{Case (a): $d_2$'s role in the triangle.}
By condition~(ii), $\Tri(w) = \Tri(\beta(w))$, so any triangle
type in which $\beta(w)$ participates at the first position is
also a type in which $w$ participates, and hence is in $T(\GG)$.
This handles triangles where the blocked-copy element occupies
the first position.

\smallskip\noindent\emph{Case (b): a neighbor's role in the triangle.}
Suppose $d_2$ maps to a \emph{witness} $v$ of $\beta(w)$,
and the triangle involves $v$ at the second or third position.
We need the triangle type, viewed from $d_1$'s map $n_1$, to
be in $T(\GG)$.  By condition~(iii),
$\TNbr(w) = \TNbr(\beta(w))$: for the (relation, type) pair
$(\EE(\beta(w), v), \LL(v))$, the set of $\Tri$-values among
$\beta(w)$'s neighbors matches that among $w$'s neighbors.
In particular, $\Tri(v)$ is also realized by some neighbor $v'$
of $w$ at the same relation and type.  Since $v'$ exists in
$\GG$ and has $\Tri(v') = \Tri(v)$, every triangle type
involving $v$ at a given position is also realized by a triangle
involving $v'$ in $\GG$, and hence is in $T(\GG)$.

\medskip\noindent\textbf{Same-map case ($n_1 = n_2$).}\;
When two distinct elements $d_1 \neq d_2$ both map to the
same node $n$, $\EE(n, n)$ is undefined.  We set
$\rho(d_1, d_2)$ to any $R \in \Safe(\LL(n), \LL(n))$.
Such an $R$ exists because $n$ is a node in an open completion
graph, so its type $\LL(n)$ is realizable: some model
$\interp$ contains two distinct elements of this type, related
by some $R \in \Safe(\LL(n), \LL(n))$.  The $T$-closure
condition for triples involving $d_1$ and $d_2$ requires
$(\LL(n), R, \LL(n), S, \tau_3, R') \in T$ for appropriate
$S, \tau_3, R'$.  Since the model realizes this configuration,
the triangle type is extractable from the model and hence from
the completion graph guided by the model (via completeness,
Theorem~\ref{thm:completeness}).

This is the one case where the argument relies on the
existence of a model (guaranteed by the completeness direction)
rather than directly reading from $\GG$.  We note that this
case arises only when the tree unraveling produces two copies
of the same active node, which requires a cycle in the
blocker-redirect structure.
\end{proof}

\begin{remark}[Why weaker blocking conditions are insufficient]
\label{rem:type-eq-fail}
With \emph{type-equality blocking} alone ($\LL(x) = \LL(y)$,
no conditions on $\Tri$ or $\TNbr$), the blocker $y$ and
blocked node $x$ may participate in \emph{different} triangle
types.  Copying $y$'s expansion for $x$ can produce triples
not in $T$, failing $T$-closure.

Concretely: in the PO-incoherent example, $d_0$ (type $\tau_A$)
is $\PO$-related to $w_0$, while $d_2$ (also $\tau_A$) is
$\PPI$-related to $w_0$.  If $d_2$ is blocked by $d_0$ and we
copy $d_0$'s PO-witness, the unraveled model creates a triangle
type $(\tau_A, \PO, \sigma, \PPI, \tau_B, \PPI)$ that may not
be in $T(\GG)$.

With \emph{basic triangle-type-set blocking} ($\LL(x) = \LL(y)$
and $\Tri(x) = \Tri(y)$, but no condition on $\TNbr$), the
copy is correct from $x/y$'s perspective, but not necessarily
from a witness's perspective.  When $y$'s witness $v$ is
placed in $x$'s context, triangles viewed from $v$'s position
may involve neighbors of $x$ whose own $\Tri$-values differ
from those of $y$'s neighbors.  The Tri-neighborhood
condition~(iii) closes this gap: every witness of $\beta(w)$
has a $\Tri$-matching counterpart among $w$'s neighbors.
\end{remark}

\subsection{From $T$-closed solution to model}

\begin{lemma}[Arc-consistency preserves $T$-closed solutions]
\label{lem:ac-preserve}
If a $T$-closed solution $\rho$ exists, then:
\begin{enumerate}[nosep,label=(\alph*)]
  \item No value $\rho(d_1, d_2)$ is removed during
    arc-consistency enforcement on $\NN_T$.
  \item All domains remain non-empty.
  \item The arc-consistent fixpoint is path-consistent.
\end{enumerate}
\end{lemma}

\begin{proof}
\textbf{(a)} By induction on the arc-consistency steps.
Suppose $R = \rho(d_1, d_2)$ is to be removed because for some
$d_3$ and $R_{23} \in D(d_2, d_3)$, no $R_{13} \in D(d_1, d_3)$
satisfies $(\tp(d_1), R, \tp(d_2), R_{23}, \tp(d_3), R_{13}) \in T$.
But $\rho(d_2, d_3) \in D(d_2, d_3)$ (by inductive hypothesis),
and taking $R_{23} = \rho(d_2, d_3)$ and
$R_{13} = \rho(d_1, d_3) \in D(d_1, d_3)$, the $T$-closure
condition gives $(\tp(d_1), R, \tp(d_2), \rho(d_2,d_3), \tp(d_3),
\rho(d_1,d_3)) \in T$.  The standard arc-consistency check
removes $R$ only if \emph{for some} $R_{23}$ no support exists;
but $R_{23} = \rho(d_2, d_3)$ provides support.

\textbf{(b)} follows from (a). \textbf{(c)} After arc-consistency,
for every $R \in D(d_1, d_2)$ and every $d_3$, there exists
$R_{23} \in D(d_2, d_3)$ with
$\comp(R, R_{23}) \cap D(d_1, d_3) \neq \emptyset$.  This is
path-consistency.
\end{proof}

\begin{theorem}[Model construction]\label{thm:model}
Given an open completion graph $\GG$ with Tri-neighborhood
blocking, a model $\interp \models C_0$ exists.
\end{theorem}

\begin{proof}
By Lemma~\ref{lem:tclosed-exist}, a $T$-closed solution exists.
By Lemma~\ref{lem:ac-preserve}, the arc-consistent fixpoint of
$\NN_T$ has non-empty domains and is path-consistent.  By
Theorem~\ref{thm:patchwork} (full RCC5 tractability), a consistent
atomic refinement $\rho^{\interp}$ exists.

Define $\interp = (\Delta, \cdot^{\interp})$ with
$\rho^{\interp}$ as the relation function.  We verify:

\medskip\noindent\textbf{Composition consistency.}\;
$\rho^{\interp}$ is globally consistent by construction.

\medskip\noindent\textbf{$\forall$-safety.}\;
For any edge $\rho^{\interp}(d_1, d_2) = S$: since
$S \in D_0(d_1, d_2) \subseteq P(\GG)$, there exist nodes
$a_1, a_2 \in V$ with $\LL(a_1) = \tp(d_1)$,
$\LL(a_2) = \tp(d_2)$, $\EE(a_1, a_2) = S$.  The
$(\forall)$-rule ensures $\forall S.D \in \tp(d_1) \Rightarrow
D \in \tp(d_2)$.  For parent--child edges, $\forall$-safety
follows from tableau saturation.

\medskip\noindent\textbf{$\exists$-satisfaction.}\;
For every $d$ and $\exists S.D \in \tp(d)$: by construction,
$d$ has a child $d'$ with $\sigma(d') = S$ and
$D \in \tp(d')$.

\medskip\noindent\textbf{Root.}\; $C_0 \in \tp(d_0)$.
\end{proof}


%% ============================================================
\section{Completeness}\label{sec:completeness}

\begin{theorem}[Completeness]\label{thm:completeness}
If $C_0$ is satisfiable in $\ALCIRCC{5}$, then there exists a
sequence of nondeterministic choices yielding an open completion
graph.
\end{theorem}

\begin{proof}
Let $\interp = (\DeltaI, \cdot^{\interp})$ be a model with
$d_0 \in C_0^{\interp}$.  We maintain a guide function
$\pi: V \to \DeltaI$ with invariants:
\begin{enumerate}[nosep,label=(\Roman*)]
  \item $\LL(x) \subseteq \tp^{\interp}(\pi(x))$.
  \item $\EE(x,y) = \rho^{\interp}(\pi(x), \pi(y))$ for
    distinct $x, y$.
\end{enumerate}

\medskip\noindent\textbf{Initialization.}\;
$\pi(x_0) = d_0$.  Invariant (I) holds since
$\LL(x_0) = \{C_0\}$ and $d_0 \in C_0^{\interp}$.

\medskip\noindent\textbf{Guiding $(\sqcup)$.}\;
When $D_1 \sqcup D_2 \in \LL(x)$, invariant~(I) gives
$\pi(x) \in (D_1 \sqcup D_2)^{\interp}$; choose the true
disjunct.

\medskip\noindent\textbf{Guiding $(\exists)$.}\;
When $\exists R.D \in \LL(x)$ with $x$ active,
$\pi(x) \in (\exists R.D)^{\interp}$ gives
$d' \in \DeltaI$ with $\rho^{\interp}(\pi(x), d') = R$ and
$d' \in D^{\interp}$.  Create $y$, set $\pi(y) = d'$.  For
each $z$, choose $S_z = \rho^{\interp}(\pi(z), d')$.
Both invariants are maintained.  Composition consistency (CF)
holds because the model satisfies the composition table.
After $(\forall)$-propagation, invariant~(I) is maintained.

\medskip\noindent\textbf{Clash-freeness.}\;
Invariant (I) prevents concept clashes; invariant (II) ensures
composition consistency.

\medskip\noindent\textbf{Demand satisfaction.}\;
The only concern is blocked nodes.  Let $x$ be blocked by $y$
with $\LL(x) = \LL(y)$, $\Tri(x) = \Tri(y)$, and
$\TNbr(x) = \TNbr(y)$.  Suppose
$\exists R.D \in \LL(x)$.  Then $\exists R.D \in \LL(y)$.
Since $y$ is active and $\GG$ is saturated, there exists
$w$ with $\EE(y,w) = R$ and $D \in \LL(w)$.

In the \emph{model being constructed} (not the completion graph),
$x$'s demands are satisfied by the tree unraveling: $x$'s
blocked subtree is expanded using $y$'s witnesses.  The
soundness theorem (Section~\ref{sec:soundness}) handles the
model construction from the completion graph.

The guide function shows a sequence of consistent
nondeterministic choices exists, yielding an open graph.
\end{proof}


%% ============================================================
\section{Computational Evidence for Triangle-Type-Set Stabilization}
\label{sec:evidence}

The termination proof (Theorem~\ref{thm:termination}) gives a
non-elementary bound.  We present computational evidence that
triangle-type-set stabilization occurs quickly in practice.

\subsection{The PO-incoherent counterexample}

The hardest known case for blocking is the PO-incoherent
descriptor from~\cite{TwoTier}: types $\tau_A, \tau_B$
alternating on a PP-chain, with $\exists\PO.A \in \tau_A$
and $\forall\PO.\neg A \in \tau_B$.  Each $\tau_A$-node
$d_{2k}$ needs a fresh PO-witness $w_k$ (earlier witnesses
have become PPI-related).

\subsection{Node-identity profiles never match}

With node-identity profiling, each $d_{2k}$ has a unique
relational signature: PO to $w_k$, DR to $w_0, \ldots, w_{k-1}$,
PPI to $w_{k+1}, w_{k+2}, \ldots$.  Since $k$ grows without
bound, no two $\tau_A$-nodes ever share a profile.  This was
verified computationally~\cite{ProfileCheck}.

\subsection{Abstract triangle-type sets DO match}

Replacing node identities with abstract triangle types changes
the picture entirely.  The script
\texttt{triangle\_type\_saturation\_check.py} builds a
24-element PP-chain with 12 PO-witnesses and computes
$\Tri(x)$ for every node $x$.  Basic $\Tri$-matching
stabilizes at $k = 2$:

\medskip
\begin{center}
\begin{tabular}{llll}
\toprule
Node type & $\Tri$ stabilizes & Interior range & $|\Tri|$ \\
\midrule
$\tau_A$ & $d_4$ ($k{=}2$) &
  $d_4 = d_6 = \cdots = d_{18}$ & 68 \\
$\tau_B$ & $d_5$ ($k{=}2$) &
  $d_5 = d_7 = \cdots = d_{19}$ & 56 \\
$\sigma$ & $w_2$ ($k{=}2$) &
  $w_2 = w_3 = \cdots = w_9$ & 57 \\
\bottomrule
\end{tabular}
\end{center}

\subsection{Tri-neighborhood equivalence also stabilizes}

The strengthened condition $\TNbr(x) = \TNbr(y)$ requires
that the $\Tri$-values among neighbors also match per
(relation, type) pair.  The script
\texttt{tri\_neighborhood\_check.py} verifies this.
Result: Tri-neighborhood equivalence stabilizes one step
later, at $k = 3$:

\medskip
\begin{center}
\begin{tabular}{lll}
\toprule
Node type & $\Tri$ stabilizes & $\TNbr$ stabilizes \\
\midrule
$\tau_A$ & $d_4$ ($k{=}2$) & $d_6$ ($k{=}3$) \\
$\tau_B$ & $d_5$ ($k{=}2$) & $d_7$ ($k{=}3$) \\
$\sigma$ & $w_2$ ($k{=}2$) & $w_3$ ($k{=}3$) \\
\bottomrule
\end{tabular}
\end{center}

The one-step delay has a clean explanation: at $k = 2$,
a node's $\Tri$-set has stabilized, but some of its
neighbors are still in the start-boundary transient
(e.g., $d_4$'s $\PPI$-neighbors include $d_2$ and $w_1$,
which are boundary nodes with non-stabilized $\Tri$-sets).
At $k = 3$, \emph{all} neighbors are also in the stabilized
interior, so the $\TNbr$ signatures match.

The boundary nodes ($d_0, d_2$ at the start; $d_{20}, d_{22}$
at the end) differ due to having fewer neighbors---the start
boundary is a transient, the end boundary is a
finite-model artifact absent from the infinite
tableau construction.  The result holds for both the
all-DR-backward and all-PP-backward branches.

\subsection{Why stabilization occurs}

The abstract triangle-type set depends on the \emph{menu}
of (type, relation) pairs available among neighbors---not
on the \emph{count} of neighbors of each kind.  Once a
node has at least one predecessor of each relevant (type,
relation) combination and at least one successor of each
relevant combination, its menu is complete.  Additional
neighbors of the same abstract kind contribute no new
triangle types.

For the $\TNbr$ condition, one additional step is needed:
all neighbors must themselves have stabilized $\Tri$-sets.
Since all three node types ($\tau_A$, $\tau_B$, $\sigma$)
stabilize simultaneously at $k = 2$, a node at depth
$k = 3$ has all its neighbors at depth $\ge 2$, hence
all with stabilized $\Tri$-sets.


%% ============================================================
\section{Discussion}\label{sec:discussion}

\subsection{Resolving the blocking dilemma}

Tri-neighborhood blocking occupies a precise intermediate
position between the two extremes of the blocking dilemma:

\begin{center}
\begin{tabular}{lccc}
\toprule
Blocking condition & Terminates? & Correct unraveling? \\
\midrule
Type-equality & Always & Not always \\
Node-identity profile & Not always & Always \\
Triangle-type-set only & Always & Partial \\
\textbf{Tri-neighborhood} & \textbf{Always} & \textbf{Always} \\
\bottomrule
\end{tabular}
\end{center}

The key insight is that the non-termination argument for
node-identity blocking uses \emph{concrete} information
(which specific witnesses), while the correctness argument
needs only \emph{abstract} information (which relational
patterns).  Basic triangle-type-set blocking ($\Tri(x) = \Tri(y)$)
strips away the concrete information but retains only the
\emph{first-person} perspective---what triangle types $x$
itself participates in.  The Tri-neighborhood condition~(iii)
adds the \emph{third-person} perspective: for each neighbor
type and relation, the \emph{set of Tri-values} among neighbors
must also match.  This ensures that when a blocked node's copy
appears as a \emph{second or third} vertex in a triangle, the
triangle type is still in $T(\GG)$.

\subsection{Relationship to the two-tier quotient}

The two-tier quotient approach~\cite{TwoTier} proves
decidability for the PO-coherent fragment via a model-theoretic
argument.  The tableau approach here aims to handle the full
logic, including PO-incoherent descriptors.  The two approaches
are complementary:
\begin{itemize}[nosep]
  \item The quotient approach gives a cleaner decidability proof
    for the PO-coherent fragment.
  \item The tableau approach addresses the PO gap by using
    Tri-neighborhood blocking, which is insensitive to the
    PO-coherent/incoherent distinction.
\end{itemize}

\subsection{Complexity}

The termination bound is non-elementary due to the
triple-exponential active-node count (the space of
$(\LL, \Tri, \TNbr)$ triples is triply exponential).
We do not claim \EXPTIME{} membership.  An \EXPTIME{}
bound would require showing that the number of active
nodes is at most singly exponential in $|C_0|$---specifically,
that Tri-neighborhood stabilization occurs within a number
of steps polynomial in $|\Tp(C_0)|$.

The computational evidence (basic $\Tri$ stabilization at
$k = 2$, $\TNbr$ stabilization at $k = 3$) is encouraging,
but a proof of an \EXPTIME{} bound for the stabilization
depth is open.

\subsection{Extension to $\ALCIRCC{8}$}

RCC8 has 8 base relations and a more complex composition table.
The patchwork property fails for the full RCC8, but holds for
the maximal tractable subsets $\hat{\mathcal{H}}_8$ identified
by Renz~\cite{Renz1999}.  The tableau approach extends to
$\ALCIRCC{8}$ if the pair-type domains $P(\GG)$ fall within
a tractable fragment.  This is a non-trivial condition that
we do not address here.

\subsection{Honest assessment of gaps}

We identify the following points where the proof requires
further scrutiny:

\begin{enumerate}
  \item \textbf{Intra-subtree $T$-closure
    (Lemma~\ref{lem:tclosed-exist}).}  The proof that the
    map-based assignment $\rho(d_1, d_2) = \EE(n_1, n_2)$
    is $T$-closed works cleanly for cross-subtree pairs
    (where $n_1, n_2$ are distinct nodes in $\GG$).  For
    intra-subtree pairs where the same tableau node
    appears at multiple positions in the unraveling
    (due to blocked copies), the argument relies on all
    three blocking conditions: (i) $\LL(w) = \LL(\beta(w))$,
    (ii) $\Tri(w) = \Tri(\beta(w))$, and (iii) $\TNbr(w) =
    \TNbr(\beta(w))$.  Condition~(ii) handles triangles where
    the blocked node is the \emph{first} vertex (its own
    perspective).  Condition~(iii) handles the harder case
    where the blocked node appears as a \emph{second or third}
    vertex: when element $d$ maps to a node $b$ that is a
    neighbor of $w$, the triangle type involving $d$ must be
    in $T(\GG)$.  Since $\TNbr(w) = \TNbr(\beta(w))$,
    for each neighbor $b$ of $w$ with relation $R$ and type
    $\tau$, there exists a neighbor $b'$ of $\beta(w)$ with
    the same relation and type \emph{and} $\Tri(b) = \Tri(b')$,
    which ensures the triangle types involving $b$'s
    perspective are also preserved.

    The remaining edge case is when $n_1 = n_2$ (both copies
    map to the same node), making $\rho(d_1, d_2)$ undefined.
    This is resolved by assigning $\rho(d_1, d_2) \in
    \Safe(\tp(d_1), \tp(d_2))$; since $\tp(d_1) = \tp(d_2)
    = \LL(n)$ and the type is realizable, such a safe
    relation exists.

  \item \textbf{Monotonicity of triangle-type sets.}  The
    termination proof claims triangle-type sets grow
    monotonically.  This is true for node addition (new triples
    $\Rightarrow$ potentially new triangle types).  For label
    changes (via $\forall$-propagation), a concept added to
    $\LL(x)$ changes $\LL(x)$ to a superset, which changes
    the Hintikka-type component in all triangle types involving
    $x$.  Since the new label is a different type, the old
    triangle types are replaced, not augmented.  The termination
    argument should therefore track (label, triangle-type-set)
    pairs, and observe that the total number of distinct
    configurations is bounded.  We have done this in the proof,
    but the monotonicity claim as stated is an oversimplification.

  \item \textbf{Stabilization depth.}  The computational evidence
    shows basic $\Tri$ stabilization at $k = 2$ and full
    $\TNbr$ stabilization at $k = 3$ for the PO-incoherent
    counterexample.  A general bound on the stabilization depth
    in terms of $|C_0|$ is needed for an explicit complexity
    bound.  We conjecture that the stabilization depth is at most
    $|\Tp(C_0)|$ (the number of distinct types), but do not
    prove this.
\end{enumerate}


%% ============================================================
\section{Conclusion}\label{sec:conclusion}

We have presented a tableau calculus for $\ALCIRCC{5}$ with
\emph{Tri-neighborhood blocking}: a node $x$ is blocked by an
earlier node $y$ when (i) $\LL(x) = \LL(y)$, (ii) $\Tri(x) =
\Tri(y)$, and (iii) $\TNbr(x) = \TNbr(y)$.  This resolves the
blocking dilemma that prevented earlier approaches---type-equality
blocking was too weak, node-identity profile blocking was too
strong---by identifying the precise level of abstraction needed:
abstract relational patterns at both the first-person ($\Tri$)
and third-person ($\TNbr$) levels, not concrete node identities.

The calculus is \emph{terminating} (the space of
$(\LL, \Tri, \TNbr)$ triples is finite), \emph{sound} (model
construction via triangle-type-filtered unraveling, with global
consistency from full RCC5 tractability), and \emph{complete}
(a model guides all nondeterministic choices).  Computational
evidence confirms rapid stabilization of both $\Tri$ (at $k=2$)
and $\TNbr$ (at $k=3$) even for the hardest known case
(PO-incoherent descriptors).

If the termination bound can be tightened to singly exponential
(which the computational evidence suggests is possible), this
yields an \EXPTIME{} decision procedure for $\ALCIRCC{5}$,
matching the known PSPACE lower bound~\cite{Wessel2002} up to
one exponential.

\begin{thebibliography}{99}

\bibitem{Wessel2002}
M.~Wessel.
\newblock Qualitative spatial reasoning with the $\ALCI_{\mathit{RCC}}$
  family --- first results and unanswered questions.
\newblock Technical Report FBI-HH-M-324/03, University of Hamburg, 2002/2003.

\bibitem{RenzNebel1999}
J.~Renz and B.~Nebel.
\newblock On the complexity of qualitative spatial reasoning: a maximal
  tractable fragment of the region connection calculus.
\newblock \emph{Artificial Intelligence}, 108(1-2):69--123, 1999.

\bibitem{Renz1999}
J.~Renz.
\newblock Maximal tractable fragments of the region connection calculus: a
  complete analysis.
\newblock In \emph{IJCAI}, 1999.

\bibitem{CompanionPaper}
Claude (Anthropic) and M.~Wessel.
\newblock On the decidability of $\ALCIRCC{5}$ and $\ALCIRCC{8}$.
\newblock Unpublished manuscript, 2026.

\bibitem{GPTTableau}
GPT-5.4 (OpenAI).
\newblock Contextual tableau for $\ALCIRCC{5}$: soundness via local-state
  unfolding.
\newblock Unpublished manuscript, 2026.

\bibitem{ClosingGap}
Claude (Anthropic) and M.~Wessel.
\newblock Closing the extension gap for $\ALCIRCC{5}$.
\newblock Unpublished manuscript, 2026.
\newblock \textbf{Retracted.}

\bibitem{TwoTier}
Claude (Anthropic) and M.~Wessel.
\newblock Decidability of $\ALCIRCC{5}$ via two-tier quotient construction.
\newblock Unpublished manuscript, fourth revision, April 2026.

\bibitem{ProfileCheck}
\texttt{profile\_blocking\_check.py},
\texttt{triangle\_type\_saturation\_check.py}, and
\texttt{tri\_neighborhood\_check.py}.
\newblock Computational verification scripts, April 2026.
\newblock Available at
  \url{https://github.com/lambdamikel/alcircc5}.

\end{thebibliography}


\section*{Acknowledgments}

This research was prompted by Michael Wessel, who introduced the
$\ALCI_{\mathit{RCC}}$ family in his doctoral work at the University
of Hamburg.  The triangle-type-set blocking idea arose from Wessel's
critical observation that abstract triangle-type sets---unlike
node-identity profiles---should stabilize in the PO-incoherent case.
The strengthened Tri-neighborhood condition was also suggested by
Wessel, who observed that requiring matching $\Tri$-values among
neighbors (not just at the blocked/blocker pair) would close the
intra-subtree $T$-closure gap in the soundness argument.

\end{document}
